% Style commun des lettres de motivation — Ismaila AMADOU DEH
\documentclass[11pt,a4paper]{article}
\usepackage[utf8]{inputenc}
\usepackage[T1]{fontenc}
\usepackage[english,french]{babel}
\usepackage{lmodern}
\usepackage[margin=2.2cm,top=1.8cm]{geometry}
\usepackage{xcolor}
\usepackage{parskip}
\usepackage[hidelinks]{hyperref}
\definecolor{accent}{HTML}{1F4E79}
\definecolor{gris}{HTML}{6B675F}
\pagestyle{empty}
% Coordonnées privées (e-mail, téléphone) : fichier local non versionné
\newcommand{\contact}{}
\IfFileExists{coordonnees.tex}{\input{coordonnees.tex}}{}
\newcommand{\entete}{%
  {\LARGE\bfseries\color{accent} Ismaila AMADOU DEH}\\[2pt]
  {\small\color{gris} Étudiant M2 Physique Fondamentale -- parcours Théories Quantiques, Université de Tours\\
  \contact\href{https://www.linkedin.com/in/ismaila-amadou-deh}{linkedin.com/in/ismaila-amadou-deh} \;·\;
  \href{https://github.com/IsmailaDEH/Ismaila}{github.com/IsmailaDEH/Ismaila}}\\[-4pt]
  {\color{accent}\rule{\linewidth}{0.6pt}}}
\newcommand{\destinataire}[2]{\null\hfill\begin{minipage}[t]{0.45\linewidth}#1\\[6pt]Tours, le \today\end{minipage}\par\medskip
  \textbf{Objet :} #2\par}
\newcommand{\signature}{\par\bigskip\hfill Ismaila AMADOU DEH}

\begin{document}
\entete
\destinataire{Thales Research \& Technology\\Service Recrutement\\Palaiseau}{Candidature au stage « Ingénieur en capteurs quantiques à base de centres NV dans le diamant pour la détection de signaux radiofréquences » (6 mois)}

Madame, Monsieur,

Étudiant en Master~2 Physique Fondamentale, parcours Théories Quantiques, à l'Université de Tours (Institut Denis Poisson), je vous adresse ma candidature pour ce stage de fin d'études. Les capteurs quantiques à centres NV sont précisément le domaine où je souhaite mettre en pratique ma formation, avant de poursuivre en thèse.

J'ai déjà manipulé ce système : en Master~1, j'ai réalisé des travaux pratiques de \textbf{magnétométrie par centres NV (ODMR)}, complétés par des TP de résonance paramagnétique électronique et d'effet Zeeman. J'y ai appris à relier la position et la largeur des raies de résonance au champ magnétique et aux temps de cohérence du spin, ce qui est le cœur de la détection de signaux radiofréquences.

Sur le plan théorique et numérique, j'ai développé seul un simulateur de \textbf{décohérence d'un qubit par l'équation de Lindblad} (intégrateur RK4 écrit à la main). Il reproduit les séquences de Rabi, Ramsey et d'écho, et j'en extrais $T_1$ et $T_2$ par ajustement ($T_2 = 17{,}14 \pm 0{,}09~\mu$s, en vérifiant la borne $T_2 \le 2T_1$). J'ai aussi modélisé la spectroscopie des bandes latérales d'un ion piégé (spectre du liouvillien, thermométrie rouge/bleu). Je sais donc simuler un protocole de mesure, ajuster des données expérimentales et en estimer les incertitudes. Tous ces projets sont en Python, testés automatiquement et publiés sur GitHub.

Ce que je peux apporter à votre équipe : une bonne compréhension de la physique du spin NV et de sa dynamique cohérente, l'habitude d'automatiser l'acquisition et l'analyse de données en Python, et une démarche rigoureuse (valider contre la théorie avant d'interpréter). En retour, je souhaite apprendre auprès de vous la mise en œuvre expérimentale d'un capteur NV et son application aux signaux RF.

Je serais heureux de vous présenter mes travaux lors d'un entretien. Je vous prie d'agréer, Madame, Monsieur, l'expression de mes salutations distinguées.
\signature
\end{document}
